\documentclass[12pt,a4paper]{article}
\usepackage[utf8]{inputenc}
\usepackage[T1]{fontenc}
\usepackage{amsmath,amssymb}
\usepackage[margin=2.5cm]{geometry}
\usepackage{setspace}
\usepackage{hyperref}

\onehalfspacing

\begin{document}

\begin{center}
    {\LARGE\bfseries Quantum Spectroscopy of the Autoregressive I:\\[4pt]
    Testing the Born Rule via Malus's Law in LLM Output Distributions\\[4pt]}

    \vspace{0.5em}
    {\large Franklin Silveira Baldo}\\
    \vspace{0.5em}
    March 2026
\end{center}

\begin{abstract}
\noindent
[STUB — to be written when experiment is executed.]
This paper tests whether sampling LLM outputs across simulated polarizer angles reproduces the $\cos^2\theta$ prediction of Malus's Law. The R\&G probe method treats the model's word-level binary output (transmitted/absorbed) as structurally isomorphic to photon counting. If the model's implicit physics encodes the Born rule, the output distribution across angles will follow $\cos^2\theta$. Deviations map the structure of the model's implicit quantum mechanics.
\end{abstract}

% Research agenda item: Quantum Spectroscopy Paper 1
% See EXPERIENCE.md "Current Research Agenda" → Immediate
% Protocol: see lab/EXPERIMENTS.md for execution infrastructure

\end{document}
